% 第12章 当前研究现状、遗留课题与未来方向
\section{当前研究现状、遗留课题与未来方向}\label{sec:future}

本章基于第 \ref{sec:survey} 章综述与第 \ref{sec:audit} 章审计，归纳均值回归策略研究截至 2026 年的现状、尚未解决的遗留课题，以及未来可重点关注的研究方向。所有判断均锚定于本综述证据链；凡属推测的均明确标注。

\subsection{当前研究现状（2020 年代中期）}

\textbf{经典框架仍是实践基线，学术重心已转移。} OU 过程、协整与距离法（Gatev 等 \cite{gatev2006}、Avellaneda--Lee \cite{avellaneda2010}、Vidyamurthy \cite{vidyamurthy2004}）仍是产业界配对交易与统计套利的事实标准；而学术前沿集中在四个迁移方向：

\begin{enumerate}
  \item \textbf{机器学习对统计套利的再实现}：从深度特征学习（R34 \cite{arxiv_dl_sa}）到在线学习（R37 \cite{arxiv_online_sa}）再到强化学习（R35 \cite{arxiv_rl_sa}、R38 \cite{arxiv_rl_mr}），第 \ref{sec:ml} 节所列 14 篇预印本构成这一方向的活跃谱系。现状判断：单点方法创新多、统一框架少；个别工作（如 R34）报告了严格的样本外与成本敏感性，但独立复现仍缺位（第 \ref{sec:checklist} 章）。
  \item \textbf{评测与基准的标准化}：PBO/DSR/多重检验（\cite{bailey2014pbo,bailey2014dsr,harveylz2016}）正从方法论文进入策略工程；开源平台 Qlib \cite{qlib} 提供了可复现的 AI 量化基准，但其基准聚焦截面因子，尚无对等的“均值回归专用基准”。
  \item \textbf{跨资产迁移}：配对/统计套利从美股扩散到东京（R43 \cite{arxiv_tse}）、波兰（R45 \cite{arxiv_poland_sa}）与加密货币（R40/R41 \cite{arxiv_crypto_drl,arxiv_crypto_coint}），新兴市场的流动性与成本结构差异成为新焦点。
  \item \textbf{微观结构尺度}：反转研究下探到秒级订单流层面（R44 \cite{arxiv_microstructure}），与行为金融的周/月反转（\cite{lehmann1990,jegadeesh1990}）在机制上分属不同范式。
\end{enumerate}

\textbf{综述视角}：Krauss（2017）\cite{krauss2017} 对配对交易五种方法（距离法、协整、时间序列、随机控制、机器学习）的系统综述给出结论——策略整体仍盈利但收益随市场效率提升而下降；这一判断与本综述第 \ref{sec:audit} 章的利润衰减证据一致，构成“现状 = 经典有效 + 边际递减 + 方法迁移”的三段式图景。

\textbf{评测污染成为共识问题}。White（2000）\cite{white2000} 以来确立的 Reality Check、PBO、DSR 等工具（\cite{sullivan1999,bailey2014pbo,bailey2014dsr}）使“未披露试验次数的高收益声明”在 2020 年代已不被视为可接受的标准；这是与 20 世纪文献最大的方法论代差。

\subsection{遗留课题（开放问题）}

\begin{enumerate}
  \item \textbf{利润衰减的机制归因}（开放度：高）。Gatev 距离法利润从 1962--1988 到 1989--2009 显著下降（\cite{gatev2006,dofaff2012}），但归因于套利拥挤、市场效率提升还是成本结构变化尚无定论；缺一个可分解“过度反应—流动性—交叉自相关”贡献的识别框架（\cite{lomackinlay1990} 只完成了初步分解）。
  \item \textbf{反转与动量的统一理论}（开放度：高）。同一资产在短/中/长尺度上分别呈现反转与动量（\cite{poterbasummers1988,mom2012}），但两者的统一建模仍缺；“慢动量+快反转”（R39 \cite{arxiv_slowfast}）是工程尝试而非理论解。
  \item \textbf{非平稳参数估计}（开放度：中--高）。OU 回归强度、半衰期与长期均值随样本漂移；在线/贝叶斯估计（Elliott 等 \cite{elliott2005}、R46 \cite{arxiv_neural_kalman}）提供了工程方案，但其统计性质（估计方差、置信区间）缺乏系统研究。
  \item \textbf{容量与拥挤的定价}（开放度：中）。统计套利收益的容量曲线、拥挤交易对价差均值的反馈效应未被标准化度量（\cite{avellaneda2010} 仅给出定性讨论）。
  \item \textbf{概率预测的可校准性}（开放度：中）。Kalman/HMM/深度概率输出的预测校准（calibration）研究几乎空白；第 \ref{sec:mapping} 章“概率/状态预测”维度在多数文献中仍为“无”。
  \item \textbf{微观结构反转的稳健性}（开放度：中）。秒级反转利润对做市规则、tick 大小与订单簿机制的敏感性缺乏跨市场验证（R44 \cite{arxiv_microstructure} 为单市场单机制结果）。
  \item \textbf{评测口径的统一}（开放度：中）。Qlib \cite{qlib} 覆盖截面因子基准，但均值回归领域缺少“同一数据、同一成本、同一评测”的公开擂台；跨论文不可比问题（第 \ref{sec:audit} 节）仍普遍。
  \item \textbf{合成数据与生成式模型}（开放度：中--高）。生成式模型用于金融时序的探索刚起步（如 arXiv:2501.03993 对合成组合数据的讨论），其作为均值回归信号增强与样本外校准工具的潜力尚未验证。
\end{enumerate}

\subsection{未来重点研究方向与路线图}

基于遗留课题与现有证据强度，给出按优先级排序的未来方向。每条附“研究问题—预期产出—与本文证据链的衔接”。

\begin{table}[htbp]
\centering
\caption{未来重点研究方向（按优先级）}
\label{tab:future}
\footnotesize
\begin{tabular}{@{}>{\raggedright\arraybackslash}p{1.0cm}>{\raggedright\arraybackslash}p{4.2cm}>{\raggedright\arraybackslash}p{4.4cm}>{\raggedright\arraybackslash}p{4.4cm}@{}}
\toprule
优先级 & 研究方向 & 研究问题 & 与现有证据链的衔接 \\
\midrule
P1 & 均值回归专用公开基准 & 统一数据/成本/评测口径的可复现擂台（对等 Qlib \cite{qlib}） & 直接消解第 \ref{sec:audit} 节不可比问题 \\
P1 & 概率/状态驱动自适应策略 & 贝叶斯 OU、隐状态过滤的在线策略与校准度量 & 补第 \ref{sec:mapping} 章“概率预测”维度的空白 \\
P2 & 多尺度统一建模 & 反转+动量+Regime 的单一框架与可识别机制分解 & 回应“利润衰减机制”遗留课题 \\
P2 & 跨市场迁移与元学习 & 加密→股票→期货的迁移分数与任务共享 & 扩展现有跨市场证据（\cite{arxiv_crypto_coint,arxiv_tse,arxiv_poland_sa}） \\
P3 & 评测纪律内建（AutoML） & 将 PBO/DSR 集成进策略搜索的每一轮 & 把 \cite{bailey2014pbo,bailey2014dsr} 从审计工具变为搜索约束 \\
P3 & 微观结构反转与执行耦合 & 秒级信号与最优执行/做市的一体化 & 衔接 R44 \cite{arxiv_microstructure} 与执行层研究 \\
P3 & 生成式数据增强 & 合成市场校准均值回归信号的样本外稳健性 & 前沿线索 arXiv:2501.03993 \\
\bottomrule
\end{tabular}
\end{table}

\noindent\textbf{综合判断}：均值回归研究正处于“经典实证成熟、方法迁移活跃、评测纪律正在建立”的窗口期。最高杠杆的工作不是再发现一个“新反转因子”，而是：(1) 建立统一评测基准（消解不可比问题）；(2) 给概率/状态预测补上可校准的工程实现；(3) 用机制分解回答“反转利润从哪里来、为何衰减”。三者均可直接落地为本文第 \ref{sec:reproduce} 章实验路线的延伸。

\subsection{统一量化指标体系}\label{sec:metrics}

为使上述方向可落地、可比对，给出本路线图强制使用的统一指标集（与第 \ref{sec:reproduce} 章 P0 统一评价层一致，并补充概率预测与稳健性维度）。任何实验报告缺任一 \textbf{必报} 指标即视为不可比。

\begin{table}[htbp]
\centering
\caption{统一量化指标体系（必报项）}
\label{tab:metrics}
\footnotesize
\begin{tabular}{@{}>{\raggedright\arraybackslash}p{1.4cm}>{\raggedright\arraybackslash}p{3.2cm}>{\raggedright\arraybackslash}p{3.2cm}>{\raggedright\arraybackslash}p{2.8cm}>{\raggedright\arraybackslash}p{3.2cm}@{}}
\toprule
类别 & 指标 & 定义 & 报告要求 & 与证据链的对应 \\
\midrule
收益 & 年化超额收益 & 扣基准后的年化收益 & 必报 & 消解 \cite{gatev2006,dofaff2012} 口径差 \\
收益 & 多空价差 & 多腿组合累计收益差 & 必报 & 配对/统计套利原生口径 \\
风险 & 最大回撤 & 峰值到谷底最大跌幅 & 必报 & 收益/风险平衡 \\
风险 & 偏度/峰度 & 收益分布高阶矩 & 必报 & DSR 非正态修正输入 \\
信号 & Rank IC / ICIR & 信号与未来收益秩相关 & 必报（ML 类） & 对齐 Qlib \cite{qlib} 基准口径 \\
信号 & 换手率 & 年化持仓周转 & 必报 & 成本敏感性代理 \\
显著性 & $t$ 统计量 & 零成本 OOS 夏普的 $t$ 值 & 必报 & 多重检验阈值 \cite{harveylz2016} \\
显著性 & 缩水夏普 DSR & 校正试验次数与非正态 & 必报 & \cite{bailey2014dsr} \\
显著性 & 过拟合概率 PBO & CSCV 计算的过拟合比例 & 必报（策略搜索时） & \cite{bailey2014pbo} \\
成本 & 三档收益 & 零成本/含佣金/含冲击 & 必报 & 成本侵蚀证据 \cite{martin2011} \\
成本 & 参与率上限 & 单笔成交量参与比例上限 & 必报（有容量声明时） & 冲击模型约束 \\
概率 & Brier score / 覆盖率 & 概率预测校准与区间覆盖 & P1 概率方向必报 & 填补第 \ref{sec:mapping} 章空白 \\
稳健性 & 迁移分数 & 同号且不显著恶化的市场比例 & P2 跨市场必报 & 扩展跨市场证据 \\
\bottomrule
\end{tabular}
\end{table}

\subsection{可执行实验方案（P1--P3）}\label{sec:experiments}

每个方向给出“数据 $\to$ 步骤 $\to$ 验证门”三段式方案。全部方案强制使用第 \ref{sec:metrics} 节指标集与统一成本模型（绝对参与率 + 佣金分档）。

\begin{enumerate}
  \item \textbf{P1a 均值回归专用公开基准}。数据：CRSP 日频（1962--2024）或 Qlib 公开数据 \cite{qlib}。步骤：(1) 固化配对池与标的筛选规则；(2) 固化形成期/持有期网格（如形成 6M/12M、持有 1M/3M/6M）；(3) 固化成本模型与成交时点（收盘信号、次日开盘成交）；(4) 接入经典基线（Gatev 距离法 \cite{gatev2006}、协整法 \cite{englegranger1987}、PCA 法 \cite{avellaneda2010}）与 ML 方法（\cite{arxiv_dl_sa,arxiv_rl_sa}）。验证门：能复现 Gatev--Do\&Faff 的利润衰减曲线形态（1962--1988 高、1989--2009 低），并输出统一指标表；否则记录偏差来源。
  \item \textbf{P1b 概率/状态驱动自适应策略}。数据：CRSP 或商品期货日频。步骤：(1) 贝叶斯 OU 参数后验（含半衰期后验区间）vs MLE 点估计对照；(2) Kalman/HMM 状态过滤与阈值信号结合；(3) 对概率输出做校准（Brier score、覆盖率）。验证门：状态驱动策略的 OOS 夏普不低于点估计基线，且 Brier/覆盖率显著更优；若概率增益不带来可交易增益，如实报告为负结果。
  \item \textbf{P2a 多尺度统一建模}。数据：多品种期货日频（对齐 Moskowitz 等 \cite{mom2012} 口径）。步骤：(1) 构建短周期反转信号、中周期动量信号、Regime 状态三通道；(2) 单一风险预算下优化组合权重（对偶梯度或分层优化）；(3) 按 Regime 条件化报告收益。验证门：合并组合的信息比率（IR）高于任一单通道且高于简单等权融合；Regime 切换处无结构性回撤。
  \item \textbf{P2b 跨市场迁移研究}。数据：美股、东京 \cite{arxiv_tse}、波兰 \cite{arxiv_poland_sa}、加密货币 \cite{arxiv_crypto_coint,arxiv_crypto_drl}。步骤：同一条管线（同一配对规则、同一信号、同一成本模型）跑四市场；输出迁移分数与分市场 OOS 指标。验证门：$\geq 3$ 个市场 OOS 同号且不显著恶化；对败北市场给出机制归因（流动性、涨跌停、费率差异）。
  \item \textbf{P3a 评测纪律内建（AutoML）}。数据：任意高候选量策略集。步骤：将 DSR/PBO 作为搜索目标或硬约束纳入自动策略搜索，记录试验次数 $M$、观察数 $T$、最小回测长度（MinTRL）。验证门：搜索输出的最终候选 PBO $<0.2$ 且 DSR 通过；若不可能，报告“该搜索空间无可信策略”。
  \item \textbf{P3b 微观结构反转与执行耦合}。数据：订单簿/tick 级数据。步骤：秒级反转信号（对齐 \cite{arxiv_microstructure}）与最优执行/做市策略联合优化；报告净执行收益与参与率。验证门：扣除执行成本的净收益为正，且对不同 tick 大小/做市规则稳健。
  \item \textbf{P3c 生成式数据增强}。数据：真实收益序列 + 合成序列（如 arXiv:2501.03993 的合成组合思路）。步骤：扩散/生成模型合成 OOS 校准集；用合成样本检验信号衰减与参数稳定性。验证门：合成样本上测得的信号衰减曲线与真实 OOS 一致（作为预注册检验工具），并如实标注合成数据的边界。
\end{enumerate}

\subsection{未来研究路线图（基于遗留课题）}\label{sec:roadmap}

路线图将第 \ref{sec:future} 章 8 项遗留课题映射为四个阶段，每阶段有明确交付物与验证门（gate）。总时长与人力为示意性安排，不构成承诺。

\begin{longtable}{@{}>{\raggedright\arraybackslash}p{1.9cm}>{\raggedright\arraybackslash}p{3.4cm}>{\raggedright\arraybackslash}p{4.6cm}>{\raggedright\arraybackslash}p{4.4cm}@{}}
\caption{未来研究路线图}\label{tab:roadmap}\\
\toprule
阶段 & 时间窗（示意） & 交付物 & 验证门（gate） \\
\midrule
\endhead
Phase 0：基建 & 0--3 个月 & 统一数据层+成本层+指标层代码库；P1a 基准包 v0 & 三档成本收益管线可复现；G1：经典基线在统一口径下跑通 \\
Phase 1：P1 三方向 & 3--9 个月 & 均值回归基准擂台 v1；贝叶斯 OU vs MLE 对照报告；机制分解初版（过度反应/流动性/交叉自相关） & G2：衰减曲线复现；G3：概率策略校准达标或负结果归档 \\
Phase 2：P2 两方向 & 9--18 个月 & 多尺度统一模型；四市场迁移报告 & G4：合并 IR 超单通道；G5：$\geq 3$ 市场同号 \\
Phase 3：P3 三方向 & 18--36 个月 & AutoML 评测纪律工具；微结构执行耦合原型；生成式校准工具 & G6：搜索输出 PBO$<0.2$；G7：执行净收益为正且稳健 \\
Phase 4：沉淀 & 24--36 个月 & 论文 2--3 篇 + 开源工具包；复核清单清零 & G8：全部 B 级复核完成；G9：基准结果可被第三方复现 \\
\bottomrule
\end{longtable}

\noindent\textbf{里程碑与失败处置}：每个 gate 未通过时，先按第 \ref{sec:reproduce} 章“失败透明”原则记录偏差来源，再决定调整参数、更换口径或归档负结果；禁止在未过 gate 的情况下推进下一阶段。路线图的输入—输出关系：Phase 0 的输出（基准包）是 Phase 1--4 的唯一可比口径；所有阶段共用第 \ref{sec:metrics} 节指标集，确保跨阶段结果可互相引用。
